% Generated draft: requires mathematical and source review.
\documentclass[11pt]{article}
\usepackage[margin=25mm]{geometry}
\usepackage{fontspec}
\setmainfont{Latin Modern Roman}
\usepackage{amsmath,amssymb,mathtools,unicode-math}
\setmathfont{Latin Modern Math}
\usepackage{enumitem,needspace,xurl,hyperref}
\setlength{\parindent}{0pt}
\setlength{\parskip}{5pt}
\setlength{\emergencystretch}{3em}
\newcommand{\N}{\mathbb{N}}
\newcommand{\Z}{\mathbb{Z}}
\newcommand{\Q}{\mathbb{Q}}
\newcommand{\R}{\mathbb{R}}
\newcommand{\C}{\mathbb{C}}
\newcommand{\F}{\mathbb{F}}
\newcommand{\A}{\mathbb{A}}
\newcommand{\PP}{\mathbb{P}}
\newcommand{\E}{\mathbb{E}}
\newcommand{\eps}{\varepsilon}
\newcommand{\dd}{\,\mathrm{d}}
\begin{document}
\section*{UnsolvedMath ID: 30001720}
\textbf{Source problem code:} catalogue entry
\subsection{Definitions and mathematical statement}
For integers \(n\geq 0\) and \(0\leq j\leq n\), define
\[
B_{j,n}=\left\{(b_1,\ldots,b_n)\in\{0,1\}^n:\sum_{r=1}^{n}b_r=j\right\}.
\]
These are the binary strings of length \(n\) and weight \(j\). This is the natural combinatorial parameter range; the supplied record does not explicitly state the parameter quantifiers or whether the boundary case \(n=0\) is intended.

For a specified field \(k\) and a finite-dimensional \(\Z\)-graded \(k\)-algebra \(A\), let \(A\text{-gmod}\) denote the category of finite-dimensional graded left \(A\)-modules with degree-preserving morphisms. Its Grothendieck group \(K_0(A\text{-gmod})\) is generated by classes \([M]\), subject to \([M]=[L]+[N]\) for each short exact sequence \(0\to L\to M\to N\to0\). With the convention \((M\langle1\rangle)_d=M_{d-1}\), grading shift makes this group a \(\Z[t,t^{-1}]\)-module through \(t[M]=[M\langle1\rangle]\). A finite-dimensional algebra is quasi-hereditary if it admits a finite chain of ideals whose successive factors are heredity ideals in the corresponding quotient algebras. Here a heredity ideal \(I\) of an algebra \(B\) satisfies \(I^2=I\), is projective as a left \(B\)-module, and satisfies \(I\operatorname{rad}(B)I=0\). The source does not specify an additional compatibility condition between this structure and the grading.

\textbf{Source question, with unresolved conventions retained.} Write \(U_q(\mathfrak{sl}_2)\) for the quantum enveloping algebra named in the source, and \(V_1\) for its vector representation. Does there exist a family of finite-dimensional quasi-hereditary graded algebras \(A_{j,n}\), for the intended parameters \(n,j\), giving a categorification with
\[
K_0(A_{j,n}\text{-gmod})\cong V_1^{\otimes n}(j),
\]
and with simple modules indexed by \(B_{j,n}\)? For graded simple modules, such indexing would normally mean indexing isomorphism classes \emph{up to grading shift}, or choosing one graded lift for each such class; the source does not state which meaning is intended.

The displayed isomorphism requires a specified coefficient ring, any necessary scalar extension, and a relation between grading shift and \(q\). For example, under a generic-parameter convention, a comparison with a \(\Q(q)\)-vector space would require an expression of the form
\[
K_0(A_{j,n}\text{-gmod})\otimes_{\Z[t,t^{-1}]}\Q(q),
\]
with the map from \(\Z[t,t^{-1}]\) to \(\Q(q)\) explicitly chosen. This is an illustration of missing data, not an additional hypothesis attributed to the source.

The notation \(V_1^{\otimes n}(j)\) is not defined in the supplied record. If it denotes the subspace spanned by tensor basis vectors indexed by \(B_{j,n}\), then, under the convention that \(V_1\) has basis vectors of weights \(1\) and \(-1\), it is the weight-\(n-2j\) subspace. Such a subspace is generally not a \(U_q(\mathfrak{sl}_2)\)-subrepresentation: the raising and lowering operators move between weight spaces. On that interpretation, categorification of the quantum-group action would require categories for the different \(j\) together with suitable functors and relations; these are not specified. Consequently, the supplied assertion cannot yet be read as a fully specified isomorphism of quantum-group representations.

The supplied dated curation reports a concrete quasi-hereditary categorification for tensor powers of the quantum \(\mathfrak{sl}_2\) vector representation, while stating that a broader extracted question is not exhausted by that construction. It does not identify the additional conditions distinguishing that broader target. This draft therefore retains the existence question and its formulation issues without asserting a new result or a verified remaining open problem.
\subsection{Short English statement}
The source asks for finite-dimensional graded quasi-hereditary algebras whose Grothendieck groups model the indicated parts of tensor powers of the quantum \(\mathfrak{sl}_2\) vector representation, with simple modules corresponding to binary strings having \(j\) ones. The meaning of the indicated parts, coefficient conventions, categorical action, and remaining target require clarification. The supplied status is \texttt{partially\_solved}; this draft does not independently verify that assessment.
\subsection{Sources}
\begin{itemize}
\item Local UnsolvedMath registry, supplied complete source record: permanent ID \(30001720\), exact source problem code \texttt{catalogue entry}, title ``Quasi-Hereditary Categorifications of \(U_q(\mathfrak{sl}_2)\) Tensor Powers.'' The record belongs to the Oberwolfach Reports Open Problems collection and gives proposed year \(2011\). No individual proposer is identified.
\item Supplied report citation: \emph{Representation Theory of Quivers and Finite Dimensional Algebras} (2011), Oberwolfach Reports, pp.\ 523--608, DOI \texttt{10.4171/owr/2011/10}. Supplied source URL and literature-source link, labeled ``Oberwolfach Report 10/2011'': \url{https://doi.org/10.4171/owr/2011/10}. No problem-specific page or contributor attribution is supplied.
\item Supplied literature assessment, checked \texttt{2026-08-21}: status \texttt{partially\_solved}. The curation states that the report gives a concrete quasi-hereditary categorification for tensor powers of the quantum \(\mathfrak{sl}_2\) vector representation, but that a broader extracted categorification question is not exhausted by this construction. This is supplied dated curation, not an independently verified literature assessment; the record itself requests independent verification of time-sensitive status and cited claims.
\item The supplied statement status is \texttt{corrected\_verified}, accompanied by a claim of manual semantic, typographic, and LaTeX review of the cleanup. This records the supplied editorial assessment and does not establish mathematical correctness or independent source verification for this draft.
\end{itemize}
\end{document}
